\section{Lattice Parton Physics with LaMET}
\label{sec:lattice}
%------------------------------------------------------------------------------

Lattice gauge theory simulates continuum QCD in imaginary time on a discretized 4D Euclidean lattice. The method is characterized by the finite lattice spacing $a$ and volume $L_1\times L_2\times L_3\times T$, and input parameters such as the strong coupling and quark masses. To calculate physical quantities, one usually expects to take the continuum ($a\to0$) and infinite volume $L_i, T\to\infty$ limits, as well as tuning the quark masses so that observables such as the pion mass $m_\pi$ agrees with the physical value of $\sim 140$ MeV. There are different methods to implement the fermions on the lattice~\cite{Rothe:1992nt}, which leads to different properties of the lattice action such as chiral symmetry breaking for Wilson fermions. In the lattice calculation of hadron matrix elements, the initial and final states are generated by acting the source and sink interpolation operators on the vacuum, and the ground-state contributions are filtered out by propagating over a sufficiently large Euclidean time. A boosted hadron state can be obtained by inserting momentum into the source and sink operators through Fourier transform in the 3D spatial coordinates.

The lattice QCD calculations of parton physics using LaMET started with the exploratory studies on the simplest PDFs and the gluon helicity~\cite{Lin:2014zya,Alexandrou:2015rja,Yang:2016plb}, which yielded fairly encouraging results, demonstrating that LaMET is a viable approach.
%
In subsequent studies, more attention has been paid to the systematics,
including establishing a proper renormalization and matching procedure, simulating at the physical pion mass, removing the excited-state contamination, etc. Such studies have greatly improved the precision of the calculations, with the latest results exhibiting a reasonable agreement with phenomenological PDFs~\cite{Alexandrou:2018pbm,Lin:2018qky,Alexandrou:2018eet,Izubuchi:2019lyk,Gao:2020ito}. In the meantime, explorations have also been made on similar large momentum data using the coordinate-space factorization methods including the pseudo-PDF~\cite{Orginos:2017kos,Joo:2019bzr,Joo:2019jct,Joo:2020spy} and current-current correlation~\cite{Bali:2019dqc,Sufian:2019bol,Sufian:2020vzb}.
%
Nevertheless, dedicated large-scale efforts with the state-of-art resources are yet to be seen. Lattice parton physics with LaMET is just at its dawn. With EIC
in the US going forward, a new era of lattice calculations is to come.

In this section, we summarize the current status of lattice calculations using LaMET and discuss future prospects.
We will begin with a general discussion on what kind of lattice setups are best suited for
LaMET calculations, and then briefly summarize relevant lattice techniques that facilitate such
calculations. After that, we review the lattice calculations that have been carried out so far and
point out future improvements. A nice complementary discussion about
lattice calculations has been made in~\cite{Cichy:2018mum}. Other reviews that summarize the recent developments in the lattice calculation of PDFs can be found in~\cite{Monahan:2018euv,Zhao:2020vll,Constantinou:2020pek}.

%------------------------------------------------------------------------------
\subsection{Special Considerations for Lattice Calculations}
\label{sec:lattice-lattice}
%------------------------------------------------------------------------------

In this subsection, we discuss the challenges for lattice calculations in LaMET, and estimate the required lattice
requirements by taking the collinear PDFs as an example.

%------------------------------------------------------------------------------
\subsubsection{Challenges due to large momentum}
\label{sec:lattice-lattice-challenges}
%------------------------------------------------------------------------------

In addition to common challenges with other lattice calculations, such as taking the continuum and infinite volume limits, simulating at or extrapolating to the physical pion mass, etc., LaMET applications
require generating large-momentum hadron on lattice. For LaMET expansion,
$1/(xP^z)$ is the expansion parameter, and for the coordinate-space factorizations, large
quasi-light-cone distance $\lambda$ requires even bigger hadron momentum.
However, realizing this faces a number of practical challenges.
First, it has been difficult to generate large-momentum hadron states on the lattice, until the technique of momentum smearing~\cite{Bali:2016lva} was proposed.
The conventional smearing method in coordinate space is designed to increase the overlap with ground state hadron at rest. Thus, it is not surprising that such a smearing is not efficient when the hadron has a large momentum.
The momentum smearing technique introduces an extra phase factor $e^{i\vec{k}\cdot \vec{z}}$ to the quark field,
such that it is peaked at nonzero momentum $\vec{k}$ in Fourier space, as shown in~\fig{lattice-momentum-smearing}.
In this way, the overlap with high momentum state is vastly increased after Euclidean time evolution. Recently, the momentum smearing technique has been incorporated into the framework of distillation~\cite{Egerer:2020hnc} to improve the extraction of ground-state energy and matrix elements at momentum $\lesssim$ 3 GeV.
%
Although there are other proposed methods to generate large momentum~\cite{Wu:2018tvt}, the momentum smearing has become a standard technique in LaMET applications.

%%%%%%%%%%%%%%%%%%%%%%%%%%%%%%%%%%%%%%%%%%%%%%%%%%%%%%%%%%%%%%%%%%%%%%%%%%%
\begin{figure}
	\includegraphics[width=0.9\linewidth]{images/lattice-momentum_smearing_fig15.pdf}
	\caption{
		Conventional smearing (left) versus momentum smearing (right)~\cite{Bali:2016lva}:
		%
		Conventional smearing has small overlap with high momentum state.
		%
		Momentum smearing shifts momentum to peak at nonzero value in momentum space.
	}
	\label{fig:lattice-momentum-smearing}
\end{figure}
%%%%%%%%%%%%%%%%%%%%%%%%%%%%%%%%%%%%%%%%%%%%%%%%%%%%%%%%%%%%%%%%%%%%%%%%%%%


Second, the proton size is frame-dependent and changes with its momentum.
%
In the proton's rest frame, simulating its structure requires that the lattice spacing be much smaller than the QCD confinement scale, i.e. $a\ll\Lambda_{\rm QCD}^{-1}$.
%
When the proton is moving fast, it undergoes Lorentz contraction by a boost factor $\gamma$ in the momentum direction, thus a finer lattice spacing $a\ll(\gamma \Lambda_{\rm QCD})^{-1}$ is needed.
%
If $a\leq 0.2$ fm is the minimum requirement to investigate a static proton, one will need at least $a\leq 0.04$ fm to have the same resolution for a proton at 5 GeV.
%
A smaller lattice spacing is difficult to achieve with current computing resources,
for it suffers from the well-known critical slowing down problem, i.e., the auto-correlation times of observables such as the topological charge increase when approaching the continuum limit~\cite{Schaefer:2010hu}, which can be much longer than the Monte-Carlo simulation times.
%
A lattice with open (Neumann) boundary condition on gauge fields in the Euclidean time direction~\cite{Luscher:2011kk}, which allows topological charge to flow in and out at boundaries of time, may overcome this problem.


Third, the gaps between the ground state and the excited state energies become smaller because of the time dilation effect.
%
In the proton's rest frame, the excited state contamination exponentially decays with the mass gap $\Delta M$ and evolution time $\tau$ in the form of $e^{-\Delta M \tau}$.
%
In the boosted frame, the mass gap $\Delta M$ in the decay factor is replaced by the energy gap $\Delta E \sim \Delta M/\gamma$, and the decay changes like $e^{-\Delta M \tau} \to e^{-\Delta E \tau}=e^{-\Delta M \tau / \gamma}$ under Euclidean time evolution.
%
Therefore, with a boosted state, a longer time evolution (source-sink separation) is needed.
%
For example, if a source-sink separation of 1 fm is needed to separate the excited state of proton with 2 GeV momentum, a proton with 5 GeV momentum will require a source-sink separation of 2.5 fm.
%
Even if the two-state fit technique is used, a longer time evolution is still required so that only the ground and first excited states dominate.

Last but not the least, lattice calculation requires $P^z\ll 1/a$, so that discretization effects of ${\cal O}((aP^z)^n)$ are under control. Therefore, one has to go to smaller lattice spacing in order to reach larger momentum. The quantification of ${\cal O}((aP^z)^n)$ effects alone in LaMET calculations has not been done, as all discretization errors are treated on equal footing in continuum extrapolation.

To summarize, to achieve a precision calculation of boosted hadron structure on lattice, a fine lattice spacing (at least in the longitudinal direction) and a large box size in the time direction are essential, which of course will also require control over the signal-to-noise ratio at large Euclidean times.

%------------------------------------------------------------------------------
\subsubsection{Considerations for lattice setup}
\label{sec:lattice-lattice-special}
%------------------------------------------------------------------------------


In practical calculations, a correlation function is first obtained on lattice in coordinate space, and then Fourier transformed to momentum space with the phase factor $e^{i\lambda x}$ where $\lambda=zP^z$.
Therefore, the smallest $x$ one can reach can be roughly estimated from the largest $\lambda$ as $x\sim 1/\lambda$.
%
However, a more stringent constraint comes from requiring that the higher-twist contribution ${\cal O}(\Lambda_{\rm QCD}^2/(xP^z)^2)$ be small so that the factorization is still valid, which implies $x\gg \Lambda_{\rm QCD}/P^z$.
%
This also provides a rough estimate for the largest attainable $x$ ($x\ll 1-\Lambda_{\rm QCD}/P^z$) since the momentum fraction carried by other partons is $\sim (1-x)$ which should also be bounded from below by the above estimate.


For the current state-of-the-art simulations, the lattice spacing can reach 0.04 fm~\cite{Fan:2020nzz,Gao:2020ito}, which implies $P^z_{\rm max}\sim 5$ GeV and the effective resolution in longitudinal direction is about $\gamma a\sim 0.2$ fm.
%
Thus the valid $x$ region that can be extracted from lattice is roughly 0.1 to 0.9.
%
On the other hand, to avoid finite volume effects, it is believed that $m_\pi L\gtrsim 4$.
%
For physical pion mass, the box size in spatial direction $L$ should be at least 6 fm, which means the box size is 150 lattice spacing. So far, the largest box size in LaMET calculations is 5.8 fm~\cite{Lin:2018qky}.
%
As discussed in Sec.~\ref{sec:lattice-lattice-challenges}, the source-sink separation of 2.5 fm is needed for $P^z=5$ GeV.
%
So the box size in time direction $T$ does not need to be particularly longer than $L$, and $T=L$ is sufficient in this lattice setup.
%
In summary, with $a=0.04$ fm at physical pion mass, one need a
$L^3\times T=150^3\times 150$ lattice to reliably extract $0.1<x<0.9$ region, which
could be possible in an exa-scale computer.

There are potential tricks to reduce the computational cost.
%
First, the required source-sink separation can be shorter if one uses a multi-state instead of two-state fit with enough statistics.
%
However, since the number of fitting parameters in $n$-state fit grows as $n^2$, such a fitting will become infeasible for too large $n$.
%
Second, note that the resolution required for transverse proton structures is not affected by the Lorentz boost, one may use a coarse lattice in the transverse directions, $a_\perp=0.1$ fm.
%
The required box size is then $L_\parallel\times L_\perp^2\times T=150\times 60^2\times 150$.
%
This asymmetric lattice can greatly reduce the resources needed for large momentum since the transverse box size is fixed.
%
However, generating configurations and renormalization on such a lattice might bring new problems and should be further studied.

In the near future, exascale supercomputers may help to reach higher momentum, as large as 5 GeV for the proton,
and improve the precision of LaMET calculations. Further theoretical developments and new ideas on the technique and algorithms are also needed to overcome the simulation difficulties.


%------------------------------------------------------------------------------
\subsection{Non-Singlet PDFs}
\label{sec:lattice-pdf}
%------------------------------------------------------------------------------

In this subsection, we review current status of lattice calculations of flavor non-singlet (isovector) PDFs in the proton and
pion.
%
The non-singlet case has the advantage that the mixing with gluons as well as the lattice calculation of disconnected diagrams can be avoided, thus greatly reducing the computational challenge. It is the most extensively studied parton observable with LaMET so far.

%------------------------------------------------------------------------------
\subsubsection{Proton}
\label{sec:lattice-pdf-proton}
%------------------------------------------------------------------------------

The pioneering lattice studies for the isovector quark PDF in the proton were carried out in~\cite{Lin:2014zya,Alexandrou:2015rja}.
These are proof-of-principle studies as the renormalization of quasi-PDFs was not well understood at that time. Nevertheless, their results encouraged the follow-up theoretical works on LaMET, including a proper renormalization and matching suitable for lattice implementations.

Certain lattice artifacts have also been studied. For example, although there is no power-divergent mixing for the quasi-PDF operators on the lattice, additional operator mixings that are not seen in the continuum can still occur if a non-chiral lattice fermion such as the Wilson-type fermion is used.
%
In~\cite{Constantinou:2017sej,Chen:2017mie} it was shown that at ${\cal O}(a^0)$ the operator for the unpolarized quark quasi-PDF, $O_{\gamma^z}(z)$, can mix with the scalar operator $O_{\bf 1}(z)$, whereas $O_{\gamma^t}(z)$ does not.
%
To reduce the systematic uncertainty from such mixing, $\Gamma=\gamma^t$ has been used since then for lattice calculations of the unpolarized quark PDF, e.g. in~\cite{Alexandrou:2017huk,Green:2017xeu,Chen:2017mzz}. Similarly, for helicity and transversity cases, one should choose $\Gamma=\gamma^5\gamma^z$ and $\Gamma=i\sigma^{z\perp}=\gamma^\perp\gamma^z$, respectively, in order to avoid the mixing.
%
It should be noted that at ${\cal O}(a)$ all $\tilde{O}_\Gamma(z)$'s can mix with others~\cite{Chen:2017mie}. Nevertheless, a fine lattice spacing can reduce these effects.

In~\cite{Alexandrou:2017huk,Chen:2017mzz,Green:2017xeu}, the nonperturbative renormalization (NPR) of the quasi-PDFs was studied in the RI/MOM scheme~\cite{Martinelli:1994ty}.
%
This scheme has several advantages: The lattice regularization scheme can be converted to $\overline{\rm MS}$ scheme through RI/MOM renormalization condition,
the computation cost is affordable,
the systematic errors can be reduced or quantified more easily, etc.
%
The works before 2018 did not include NPR and the systematics were not accurately quantified.
%
The later works have implemented the RI/MOM scheme and the corresponding perturbative matching~\cite{Constantinou:2017sej,Stewart:2017tvs,Liu:2018uuj}.
%
The coordinate-space method is also developed in parallel in~\cite{Orginos:2017kos,Cichy:2019ebf,Joo:2019jct,Joo:2020spy,Bhat:2020ktg}.
%
%
In Figs.~\ref{fig:lattice-proton_ETMC_2018} and~\ref{fig:lattice-proton_LP3_2018}, we select some most recent lattice results.
%
ETMC published the proton unpolarized, helicity and transversity PDFs with $P^z=1.4$ GeV at physical pion mass~\cite{Alexandrou:2018pbm,Alexandrou:2018eet}, and ${\rm LP}^3$ published the proton helicity PDF with unprecedented momentum $P^z=3.0$ GeV at physical pion mass~\cite{Lin:2018qky}. Recently, calculations on fine lattices~\cite{Fan:2020nzz,Alexandrou:2020qtt} and an extrapolation to the continuum limit~\cite{Alexandrou:2020qtt} have become available. The finite volume effects, which was first studied in a model~\cite{Briceno:2018lfj}, have also been investigated on the lattice lately in~\cite{Lin:2019ocg}, where no sizeable volume-dependence was observed at $P^z=1.3$ and $2.6$ GeV.

The PDFs extracted from LaMET can be useful for phenomenology by providing input in kinematic regions that are difficult to measure in experiments.
%
It has attracted attention from global fit community~\cite{Lin:2017snn,Hobbs:2019gob,Lin:2020rut,Bringewatt:2020ixn}.
%
For example, it has been found that in the large-$x$ region of unpolarized PDF the lattice result will lead to significant improvement on global fit result if it reaches an accuracy of about $10\%$~\cite{Lin:2017snn}.
%
The sea quark asymmetry~\cite{Geesaman:2018ixo} is also possible to be investigated now directly on lattice.
%
For the transversity PDF, due to the difficulty of measurement in experiment, lattice results can already have impact on improving global fit and even making predictions. In addition to the isovector cases, calculations of the strange and charm unpolarized distributions~\cite{Zhang:2020dkn}, as well as the flavor separation of light quarks in the helicity PDF~\cite{Alexandrou:2020uyt}, have also been carried out recently.
%
From early exploratory results showing qualitative behavior of PDFs to the latest results which are comparable with global fits, it has come a long way in developing new techniques (momentum smearing, renormalization, matching, etc.) and the computation resources have been steadily increased over time.
%
The systematic uncertainties in the lattice calculation of PDFs have been thorougly investigated by the ETMC~\cite{Alexandrou:2019lfo}.
Further studies on systematics such as the discretization effects and finite volume effects on various lattice ensembles are still necessary.
%
In the future, lattice QCD is expected to make a significant impact on nuclon structure.


To conclude this subsection, we would like to mention that there are also lattice studies of the isovector PDF of other baryons, $\Delta^+$ to be more concrete, using LaMET~\cite{Chai:2020nxw}.


%%%%%%%%%%%%%%%%%%%%%%%%%%%%%%%%%%%%%%%%%%%%%%%%%%%%%%%%%%%%%%%%%%%%%%%%%%%
\begin{figure}
	\includegraphics[width=0.45\textwidth]{images/lattice-proton_unpol_ETMC_2018_fig16a.pdf}
	\includegraphics[width=0.45\textwidth]{images/lattice-proton_transversity_ETMC_2018_fig16b.pdf}
	\caption{
		Proton isovector quark PDF~\cite{Alexandrou:2018pbm,Alexandrou:2018eet}:
		%
		The unpolarized PDF with $P^z$ from 0.82 to 1.4 GeV and the transversity PDF with $P^z=1.4$ GeV are in upper and lower figures.
		%
		CJ15~\cite{Accardi:2016qay}, ABMP16~\cite{Alekhin:2017kpj}, and NNPDF3.1~\cite{Ball:2017nwa} are global fits.
		%
		SIDIS is global fit and SIDIS$+g_T^{lattice}$ is global fit with lattice constraint on tensor charge $g_T^{lattice}$~\cite{Lin:2017stx}.
	}
	\label{fig:lattice-proton_ETMC_2018}
\end{figure}
%%%%%%%%%%%%%%%%%%%%%%%%%%%%%%%%%%%%%%%%%%%%%%%%%%%%%%%%%%%%%%%%%%%%%%%%%%%

%%%%%%%%%%%%%%%%%%%%%%%%%%%%%%%%%%%%%%%%%%%%%%%%%%%%%%%%%%%%%%%%%%%%%%%%%%%
\begin{figure}
	\includegraphics[width=0.45\textwidth]{images/lattice-proton_helicity_LP3_2018_fig17.pdf}
	\caption{
		The proton isovector quark helicity PDF ($P^z=3.0$ GeV)~\cite{Lin:2018qky} with red band for statistic error and grey band for statistic and systematic errors.
		%
		NNPDF1.1pol~\cite{Nocera:2014gqa} and JAM17~\cite{Ethier:2017zbq} are global fits.
	}
	\label{fig:lattice-proton_LP3_2018}
\end{figure}
%%%%%%%%%%%%%%%%%%%%%%%%%%%%%%%%%%%%%%%%%%%%%%%%%%%%%%%%%%%%%%%%%%%%%%%%%%%


%------------------------------------------------------------------------------
\subsubsection{Pion}
\label{sec:lattice-pdf-pion}
%------------------------------------------------------------------------------

The pion valence quark distribution has been extracted from various Drell-Yan data for pion-nucleon/pion-nucleus scattering, while theoretical predictions do not yield consistent results with the experimental extraction, especially in large-$x$ region~\cite{Holt:2010vj}.
%
LaMET calculations will be able to shed valuable light on how to resolve this disagreement, provided that all systematics are well under control.

In principle, calculating the pion valence PDF is easier than the proton PDF. First, the pion state is easier to produce and the quark contractions are fewer. Second, the energy gap between the first excited and ground state of the pion is much bigger than the energy gap of the proton.
%
Therefore, the excited-state contamination is easier to control.
%
The simulation was first performed in~\cite{Chen:2018fwa} with the same lattice setup and procedure used in exploratory studies of the proton PDF.
%
A more thorough study on the pion valence quark PDF was done by the lattice QCD group of BNL~\cite{Izubuchi:2019lyk}.
%
It is worth pointing out that the excited state contamination was thoroughly studied using multi-state fits, with the ground and first excited states both agreeing with the expected dispersion relations,
indicating that the excited contamination is well under control.
%
The comparison of the lattice results from quasi-PDF, pseudo-PDF and current-current correlator approach are shown in Fig.~\ref{fig:lattice-pion_valence_pseudo_2019}.
%
Note that the $\rm LP^3$~\cite{Chen:2018fwa} result was obtained using Fourier transformation and inversion of factorization formula, while other three groups used parameterization models to fit the lattice data. More dedicated effort is needed to reduce the errors, and a meaningful comparison between different operators and analysis methods should be made.


For other mesons, we would like to mention that there is a study of kaon valence quark PDF using MILC configurations~\cite{Lin:2020ssv}.

%%%%%%%%%%%%%%%%%%%%%%%%%%%%%%%%%%%%%%%%%%%%%%%%%%%%%%%%%%%%%%%%%%%%%%%%%%%
\begin{figure}
	\includegraphics[width=0.45\textwidth]{images/lattice-pion_valence_pseudo_2019_fig18.pdf}
	\caption{Pion valence quark PDFs in various approach:
		Compare the results of pesudo-PDF [Reduced pseudo-ITD~\cite{Joo:2019bzr}], quasi-PDF [quasi-PDF-1~\cite{Izubuchi:2019lyk} and quasi-PDF-2~\cite{Chen:2018fwa}], and the current-current correlator approach [LCSs~\cite{Sufian:2019bol}].
	}
	\label{fig:lattice-pion_valence_pseudo_2019}
\end{figure}
%%%%%%%%%%%%%%%%%%%%%%%%%%%%%%%%%%%%%%%%%%%%%%%%%%%%%%%%%%%%%%%%%%%%%%%%%%%

%------------------------------------------------------------------------------
\subsection{Gluon Helicity and Other Collinear Parton Properties}
\label{sec:lattice-other_distributions}
%------------------------------------------------------------------------------

In this subsection, we summarize the applications of LaMET to other collinear parton observables, including the gluon helicity, the gluon PDFs, meson DAs and GPDs.

\subsubsection{Total gluon helicity}
\label{sec:lattice-gluon-helicity}

The total gluon helicity $\Delta G$ is a key component in understanding the proton spin structure. It has been intensively explored at RHIC and will be dedicatedly pursued at EIC in the future. However, a theoretical lattice calculation of $\Delta G$ had not been possible until the proposal of LaMET.


The first such effort was made by $\chi$QCD collaboration in~\cite{Yang:2016plb}.  The calculation was carried out with valence overlap fermions on $2+1$ flavor domain-wall fermion gauge configurations, using ensembles with multiple lattice spacings and volumes including one with physical pion mass.
%
The authors simulated proton matrix elements of the free-field operator $(\vec E\times \vec A)^3$ in the Coulomb gauge at various momenta, and then converted them to the $\overline{\rm MS}$ scheme with one-loop lattice perturbation theory. The $\overline{\rm MS}$ matrix elements at each lattice momentum are shown in~\fig{lattice-gluon_helicity_2016}. Though a LaMET matching is necessary to match the results to the physical gluon helicity, the authors did not apply it due to the concern of perturbative convergence of the matching coefficient~\cite{Ji:2014lra}. Instead, as the $\overline{\rm MS}$ matrix elements show rather mild momentum dependence up to the maximum momentum $\sim$1.5 GeV, they extrapolated the results to infinite momentum, as well as physical pion mass and continuum limits, with a model motivated by chiral EFT. Their final result is $\Delta G(\mu^2 = 10\ \text{GeV}^2) = 0.251(47)(16)$, or 50(9)(3)\% of the total proton spin, which agrees with the truncated moment of $\Delta g(x)$~\cite{deFlorian:2014yva,Nocera:2014gqa} within uncertainties.


Despite such progress, one should be cautious that this calculation still needs further improvements in the future. Among others, the most important ones are simulations at larger proton momentum, performing an NPR and investigating perturbative convergence of LaMET matching and its implementation.

%%%%%%%%%%%%%%%%%%%%%%%%%%%%%%%%%%%%%%%%%%%%%%%%%%%%%%%%%%%%%%%%%%%%%%%%%%%
\begin{figure}
	\includegraphics[width=0.45\textwidth]{images/lattice-gluon_helicity_2016_fig19.pdf}
	\caption{Total gluon helicity~\cite{Yang:2016plb}:
		The results are extrapolated to the physical pion mass and continuum as a function of the proton momentum $p_3$ on all the five ensembles indicated by different colors of the data points.
		%
	}
	\label{fig:lattice-gluon_helicity_2016}
\end{figure}
%%%%%%%%%%%%%%%%%%%%%%%%%%%%%%%%%%%%%%%%%%%%%%%%%%%%%%%%%%%%%%%%%%%%%%%%%%%


%------------------------------------------------------------------------------
\subsubsection{Gluon PDF}
\label{sec:lattice-Gluon}
%------------------------------------------------------------------------------


The gluon PDF is of great interest not only for precision physics at LHC, but also for understanding the gluonic structure of the proton and nuclei---as well as the small-$x$ dynamics---at the future EIC.
%
With the recent progress on the renormalization and matching for gluon quasi-PDFs~\cite{Wang:2017eel,Wang:2017qyg,Zhang:2018diq,Li:2018tpe,Wang:2019tgg} or the coordinate-space ``pseudo distributions''~\cite{Balitsky:2019krf}, a systematic lattice calculation of the gluon PDFs can be carried out in principle.

Before the above theoretical developments, an exploratory lattice study of the proton and pion unpolarized gluon PDFs were carried out in~\cite{Fan:2018dxu}. The authors calculated quasi gluon LF correlations and compared them to the LF correlations for the gluon PDFs. Later on, based on the multiplicative renormalizability of certain choice of the quasi gluon LF correlator~\cite{Zhang:2018diq}, the same authors used the ratio scheme~\cite{Balitsky:2019krf} in coordinate space to renormalize the lattice matrix elements, and fitted the proton unpolarized gluon PDF with a simple two-parameter model ~\cite{Fan:2020cpa}. Although the results show agreement with the global analyses in the large-$x$ region, the systematics from the model-dependence of the fit remains to be quantified for a controlled calculation of the gluon PDF.

%------------------------------------------------------------------------------
\subsubsection{DA}
\label{sec:lattice-DA}
%------------------------------------------------------------------------------

According to \sec{others-da}, LaMET can be readily applied to calculating DAs, and the lattice resource needed is expected to be cheaper than that for PDFs since there is one less external state, which reduces the number of contractions for the quark propagators.
%
So far there are a few exploratory investigations on meson DAs, in particular, on pion~\cite{Zhang:2017bzy} and kaon DAs~\cite{Chen:2017gck}.
%
The lattice calculations of pion~\cite{Zhang:2017bzy} and kaon DAs~\cite{Chen:2017gck} were first explored without the NPR and the corresponding matching. Recently, the pion and kaon DAs from the RI/MOM scheme analysis are extrapolated to the continuum limit~\cite{Zhang:2020gaj}, where the authors eventually adopted a two-parameter model to fit the final result.
%
The above results are shown in~\fig{lattice-pion_DA_LP3_2017}.
%
Apart from LaMET, the current-current correlation methods~\cite{Braun:2007wv,Braun:2015axa,Detmold:2005gg} have also made much progress on the pion DA~\cite{Bali:2019dqc,Detmold:2018kwu,Detmold:2020lev}.
%

%%%%%%%%%%%%%%%%%%%%%%%%%%%%%%%%%%%%%%%%%%%%%%%%%%%%%%%%%%%%%%%%%%%%%%%%%%%
\begin{figure}
	\includegraphics[width=0.45\textwidth]{images/lattice-pion_DA_LP3_2017_fig20.pdf}
	\caption{Pion DA~\cite{Chen:2017gck}:
		Comparison of $\phi_\pi$ (Lat LaMET) to previous determinations in the literature.
		%
		Lat Mom 1 and 2 are parameterized fits to the lattice moments~\cite{Braun:2015axa};
		DSE is Dyson-Schwinger equation calculations~\cite{Chang:2013pq};
		Asymp is the asymptotic form $6x(1-x)$;
		Belle is a fit to the Belle data~\cite{Agaev:2012tm};
		LFCQM is light-front constituent quark model~\cite{deMelo:2015yxk}.
	}
	\label{fig:lattice-pion_DA_LP3_2017}
\end{figure}
%%%%%%%%%%%%%%%%%%%%%%%%%%%%%%%%%%%%%%%%%%%%%%%%%%%%%%%%%%%%%%%%%%%%%%%%%%%

%------------------------------------------------------------------------------
\subsubsection{GPD}
\label{sec:lattice-GPD}
%------------------------------------------------------------------------------

As discussed in \sec{gpo}, the global fitting of GPDs still faces challenges from their complicated kinematic dependence and limited information from the experimental observables despite the progress made~\cite{Kumericki:2016ehc,Favart:2015umi}.
%
On the other hand, previous lattice QCD method is only able to calculate the lowest few moments of the GPDs~\cite{Hagler:2009ni}, which is far from sufficient to reconstruct their full kinematic dependence.
Applying LaMET to GPD calculations will provide important information on the GPDs, especially in kinematic regions that are not accessible in currently available experiments.
%
In addition, on the lattice one can study the GPD dependence on one kinematic variable by fixing the others. All these will help differentiate commonly used
models in GPD parameterization.



Calculating the quasi-GPDs requires more resources than quasi-PDF, but does not need further techniques in principle.
%
Besides, the lattice renormalization factors for the quasi-PDFs can be used here, as has been argued in \sec{gpo}.
%
The first lattice calculation of the pion unpolarized isovector quark GPD was carried out in~\cite{Chen:2019lcm}, though the results are not yet able to differentiate models or compare to experiments.
Recently, ETMC completed the first proof-of-principle calculation of the proton unpolarized and helicity GPDs~\cite{Alexandrou:2020zbe}, as shown in \fig{gpd}, which demonstrates that it is feasible to extract the GPDs with controlled systematics on available computational resources.

%%%%%%%%%%%%%%%%%%%%%%%%%%%%%%%%%%%%%%%%%%%%%%%%%%%%%%%%%%%%%%%%%%%%%%%%%%%
\begin{figure}
	\includegraphics[width=0.45\textwidth]{images/lattice_gpd_fig21.pdf}
	\caption{Proton unpolarized isovector quark GPD~\cite{Alexandrou:2020zbe} $H(x,\xi,t)$ for $t=-0.69\ {\rm GeV^2}$ extracted from quasi-GPDs at $P_3=1.25$ GeV, which is compared to the unpolarized PDF $f_1(x)$.
	}
	\label{fig:gpd}
\end{figure}
%%%%%%%%%%%%%%%%%%%%%%%%%%%%%%%%%%%%%%%%%%%%%%%%%%%%%%%%%%%%%%%%%%%%%%%%%%%
%

\subsubsection{Higher-twist PDF}

The higher-twist PDFs probe multi-parton correlations, and their contribution at $x=0$ can shed light on the LF zero modes~\cite{Ji:2020baz}. As we discussed in \sec{gpo}, such distributions can also be calculated on the lattice with the LaMET approach.

The first attempt to calculate the isovector twist-three PDF $g_T(x)$ has been carried out by ETMC~\cite{Bhattacharya:2020cen} using the one-loop matching coefficient they computed in~\cite{Bhattacharya:2020xlt,Bhattacharya:2020jfj}. Their results show agreement with the Wandzura-Wilczeck approximation~\cite{Wandzura:1977qf}, which ignores the contribution from dynamical twist-three contributions, and the Burkhardt-Cottingham sum rule~\cite{Burkhardt:1970ti}. Nevertheless, the mixing between $g_T(x)$ and other twist-three distributions was not considered, and further study is still required for an accurate matching to the light-cone PDF.

%------------------------------------------------------------------------------
\subsection{TMDs}
\label{sec:lattice-TMD}
%------------------------------------------------------------------------------

With tremendous experimental focus on the TMDPDFs for studying 3D proton structures and gluon saturation at EIC, their first-principle calculation from lattice QCD will significantly boost this direction by providing useful nonperturbative inputs for all the phenomenological analyses.

In this subsection, we discuss the status and prospects of calculating the quasi-TMDPDF and soft function with LaMET. Besides, we note that before LaMET there had already been efforts to extract information of TMDs by studying ratios of the lattice correlators~\cite{Hagler:2009mb,Musch:2010ka,Musch:2011er,Engelhardt:2015xja,Yoon:2017qzo}, which has made a series of progress in the past decade. We begin with a brief review of them.

%------------------------------------------------------------------------------
\subsubsection{Pre-LaMET study --- ratio of lattice correlators}
%------------------------------------------------------------------------------

By employing Lorentz covariance, the $x$-moments of TMDPDFs are related to the form factors of spacelike staple-shaped gauge link operators, which can be directly simulated on the lattice. Although the lattice calculation of the soft function was not available during that time, ratios of the spin-dependent and the unpolarized matrix elements were formed to cancel it, thus providing useful information of different TMDPDFs.
For example, the time-reversal odd TMDPDFs can be studied with the staple-shaped gauge link operator in a transversely polarized proton state, thus helping understand properties related to single-spin asymmetry (SSA), which was measured experimentally at STAR~\cite{Adamczyk:2015gyk}
and COMPASS~\cite{Aghasyan:2017jop}.
%
In~\cite{Musch:2011er,Engelhardt:2015xja}, the Sivers and Boer-Mulders functions of proton and pion were studied; Other time-reversal even functions, such as the worm-gear function $g_{1T}$~\cite{Tangerman:1994eh}, were also studied~\cite{Yoon:2017qzo}.

%------------------------------------------------------------------------------
\subsubsection{Quasi-TMDPDF and Collins-Soper kernel}
%------------------------------------------------------------------------------

The lattice calculation of the quasi-TMDPDF defined in \eq{quasi_TMD} is straightforward.
The matrix element of the staple-shaped quark Wilson line operator can be simulated the same way as the quasi-PDF case, except that the geometry of the gauge-link is different, while the calculation of Wilson loop $Z_E$ is standard practice in lattice QCD. The more challenging part, however, is the renormalization of the quasi-TMDPDF and its matching to the $\overline{\rm MS}$ scheme.

Using the auxiliary field theory formalism, one can argue that staple-shaped quark Wilson line operator is also multiplicatively renormalizable~\cite{Ebert:2019tvc,Green:2020xco}. On a non-chiral lattice, it suffers from finite mixing with other quark bilinear operators, as was predicted by one-loop lattice perturbation theory~\cite{Constantinou:2019vyb}. The full mixing pattern for such operators with different Dirac matrices have been studied in the RI/MOM scheme on three quenched lattice ensembles with different spacings~\cite{Shanahan:2019zcq}, and a diagonalization of the mixing matrix is adopted to renormalize these operators. Meanwhile, the one-loop conversion factors that convert the RI/MOM matrix elements to the $\overline{\rm MS}$ scheme have been calculated in continuum perturbation theory for both the $z=0$~\cite{Constantinou:2019vyb} and $z\neq0$~\cite{Ebert:2019tvc} cases.

Although the soft function is still needed to fully determine the physical TMDPDF, the $\overline{\rm MS}$ quasi-TMDPDF can already be used to extract the Collins-Soper kernel according to \eq{TMD-CS_kernel}~\cite{Ji:2014hxa,Ebert:2018gzl,Ebert:2019okf}. Since the Collins-Soper kernel can be defined from both the bare TMDPDF and the soft function, it is independent of the external state and can be calculated in a pion which is the least expensive on the lattice. Up to mass corrections suppressed by the momentum in \eq{quasiTMDmatch}, this calculation also allows for using an unphysical valence pion mass, as long as the sea quark masses are physical.

With the method developed in~\cite{Ebert:2018gzl}, the first exploratory lattice calculation of the Collins-Soper kernel was performed in~\cite{Shanahan:2020zxr} on a quenched lattice with heavy valence pion mass $m_\pi\sim 1.2$ GeV, and the result is shown in~\fig{lattice-cs}. As one can see, the lattice prediction is robust for $0.1\ {\rm fm} < b_\perp < 0.8\ {\rm fm}$, which covers the nonperturbative region that is important for TMD evolution in global analyses. Besides, at small $b_\perp$, the perturbative calculation can serve as a calibration for estimating the systematic uncertainties, as there are power corrections of ${\cal O}(1/(P^z b_\perp))$ which can only be reduced with larger $P^z$. In~\cite{Zhang:2020dbb}, the Collins-Soper kernel has also been extracted from a pion quasi-TMD DA, where the lattice renormalization was left out, and the result is in agreement with~\cite{Shanahan:2020zxr} within errors for a wide range of $b_\perp$.
With improved lattice ensembles and systematic corrections in the future, it is promising to have a precise determination of the Collins-Soper kernel for TMD phenomenology.

\begin{figure}[htb]
	\includegraphics[width=0.48\textwidth]{images/lattice-cs_fig22.pdf}
	\caption{The Collins-Soper kernel from the first exploratory calculation on a quenched lattice~\cite{Shanahan:2020zxr}. The results are obtained by using fits to the $\MS$ unsubtracted quasi-TMDPDFs with Hermite and Bernstein polynomial bases. The solid and dashed lines are the perturbative predictions~\cite{Li:2016ctv,Vladimirov:2017ksc}, which is hit the Landau pole near $b_\perp\sim 0.25$ fm. The background shading density is proportional to a naive estimate of the power corrections $1/(b_\perp P^z) + b_\perp /L$.
	}
	\label{fig:lattice-cs}
\end{figure}

%------------------------------------------------------------------------------
\subsubsection{Soft function}
\label{sec:lattice-SF}
%------------------------------------------------------------------------------

As the remaining piece towards physical TMDPDFs, the soft function must be calculated in lattice QCD. In particular, the reduced soft function in \eq{reduced} eliminates the regulator-scheme-dependence of the off-the-light-cone quasi-TMDPDF, so its calculation alone has great physical significance. According to \secs{offlight}{others}, two methods have been proposed to calculate the off-the-light-cone soft function or reduced soft function on the lattice~\cite{Ji:2019sxk}, as we discuss in the following. One relies on simulating HQET on the lattice, while the other requires calculating a light-meson form factor of transversely-separated current products.

The latter method has been implemented in the first exploratory lattice calculation of the reduced soft function~\cite{Zhang:2020dbb}, which includes simulations of the pion form factor in two external states with opposite large momenta, as well as the pion quasi-TMD DA.
The results for the reduced soft function, which are obtained with tree-level matching and omission of lattice renormalization, are shown in \fig{lattice_soft}. As one can see, they agree with the perturbative prediction for small $b_\perp$ within errors, as expected, and start to deviate when $b_\perp$ becomes large. Since the quasi-TMD DA depends on the momentum $P^z$, the stability of results at different $P^z$ suggests the validity of \eq{form_qfac}. In the future, larger statistics and improved systematics in both lattice and perturbative matching will be necessary to achieve a precision calculation of this quanity.

\begin{figure}[htb]
	\includegraphics[width=0.48\textwidth]{images/lattice_soft_fig23.pdf}
	\caption{The reduced soft factor as a function of $b_\perp$ extracted from the light-meson form factor in \sec{tmd}~\cite{Zhang:2020dbb}. The results are obtained with quasi-TMD DAs at different pion momentum $P^z$, with perturbative matching and power corrections ignored. The dashed line is one-loop prediction in perturbation theory, which hits the Landau pole at $b_\perp\sim 0.3$ fm, and the grey band is the error by varying $\mu$ by a factor of $1/\sqrt{2}$ and $\sqrt{2}$.
	}
	\label{fig:lattice_soft}
\end{figure}



%------------------------------------------------------------------------------
